\begin{enumerate}[label=\textbf{\arabic*.},leftmargin=*]
\item Localization gives \((S/I)_f\cong S_f/IS_f\).  Since the quotient map is graded, taking degree-zero pieces gives the desired isomorphism.
\item A homogeneous prime of \(S/I\) is \(\mathfrak p/I\) for a homogeneous prime \(\mathfrak p\supseteq I\); it lies in Proj exactly when \(S_+\not\subseteq\mathfrak p\).
\item The quotient is \(k[x_0,x_1]\), hence the Proj is \(\mathbb P^1_k\).
\item On \(X\ne0\): \(z-y^2=0\). On \(Y\ne0\): \(xz-1=0\). On \(Z\ne0\): \(x-y^2=0\), using the corresponding ratio coordinates.
\item \(S\) has dimension \(n+1\).  Every minimal prime over \((F)\) has height one, so the quotient components have dimension \(n\), and their nonempty Proj components have dimension \(n-1\).
\item The sum is \((x_0,x_1)\), whose Proj quotient is \(\Proj k[x_2]\), a single \(k\)-point.
\item Since \(x_0,x_1\) are relatively prime prime elements, \((x_0)\cap(x_1)=(x_0x_1)\), defining the union of the two coordinate lines.
\item Both have support \(x_0=0\), but in the first quotient the class of \(x_0\) is nonzero with square zero on suitable charts.
\item Since \(\mathfrak m x_0\subset I\), \(x_0\in I^{\rm sat}\), so \((x_0)\subseteq I^{\rm sat}\).  As \(I\subseteq(x_0)\) and \((x_0)\) is saturated, equality follows.
\item If \(S_+^r\subseteq I\), any prime containing \(I\) contains \(S_+\), so no such prime lies in Proj.
\item Empty Proj means every homogeneous prime containing \(I\) contains \(S_+\), hence \(S_+\subseteq\sqrt I\).  Finite generation of \(S_+\) gives \(S_+^r\subseteq I\) for some \(r\).
\item The homogenization is \(x_0x_2-x_1^2\).  At infinity \(x_0=0\), hence \(x_1=0\), leaving the point \([0:0:1]\).
\item The homogenization is \(x_1x_2-x_0^2\).  Setting \(x_0=0\) gives \(x_1x_2=0\), hence the two points \([0:1:0]\) and \([0:0:1]\).
\item The cone is \(\Spec k[X,Y,Z]/(XZ-Y^2)\); the vertex is the maximal homogeneous ideal \((X,Y,Z)\).
\item Every \(A_{(f)}\) is a subring of the reduced localization \(A_f\), hence is reduced.  The standard charts cover Proj.
\item Every chart ring \(A_{(f)}\) is a subring of the domain \(A_f\), hence a domain.  Nonempty overlapping affine charts give an irreducible reduced scheme.
\item In \(A_{x_0}\), the relation \(x_0z=0\) and invertibility of \(x_0\) force \(z=0\); similarly on \(A_{x_1}\).
\item The homogeneous minimal primes are \((x_0)\) and \((x_1)\), neither containing the irrelevant ideal.  They give the two projective-line components.
\item The point ideals are \((x_1,x_2)\) and \((x_0,x_2)\).  Their intersection is \((x_2,x_0x_1)\).
\item On \(x_0\ne0\), let \(u=x_1/x_0\), \(v=x_2/x_0\).  The ideal becomes \((u,v)^2\), the first infinitesimal thickening of the origin in \(\mathbb A^2\).
\item The standard embedding gives \(k[s,t]\).  The quadratic Veronese image is the conic with coordinate ring \(k[X,Y,Z]/(XZ-Y^2)\).  VI/38 proves the corresponding closed embedding is an isomorphism onto that image.
\item If \(g\in I^{\rm sat}\), choose \(r\) with \(\mathfrak m^r g\subseteq I\).  Then \(x_i^r g\in I\), so after inverting \(x_i\), \(g\in IS_{x_i}\).  The reverse containment is immediate.
\item The scheme-theoretic intersection with infinity is defined by adding \((x_0)\) to the saturated homogeneous ideal of the closure.  Merely setting \(x_0=0\) in unsaturated generators can preserve spurious components.
\item On a chart where some \(x_j\ne0\), the ratio \(u=x_0/x_j\) satisfies \(u^2=0\) modulo \((x_0^2)\), whereas modulo \((x_0)\) it vanishes.  Thus the radical removes a nilpotent normal direction.
\end{enumerate}
